\chapter{Physics: from fluid dynamics to plasmas, and what rusty-SUNDIALS brings}\label{ch:physics}
\chaptermark{Physics: fluids, plasmas, rusty-SUNDIALS}

The theorem of the previous chapters is a statement about exact potentials and infinite point configurations.
Physics never offers either. It offers finite systems at finite temperature, with interaction laws that are only
approximately of the ideal form, confined by walls, bowls or electrodes, and evolving by a dynamics that may or
may not seek the minimum of anything. This chapter asks where, in real systems, the triangular lattice appears,
which of those systems are actually inside the class of the theorem, and which only resemble it.

The primary domain is \emph{fluid dynamics}: two-dimensional incompressible flow and its point vortices
(Section~\ref{sec:ch4-vortices}), and colloidal monolayers at fluid interfaces, whose overdamped motion through a
viscous liquid is exactly the gradient-flow model used in Chapter~3 (Section~\ref{sec:ch4-colloids}). We then visit
plasmas (Section~\ref{sec:ch4-plasmas}), summarise four laboratory protocols (Section~\ref{sec:ch4-lab}), and ask
what the solver library rusty-SUNDIALS~\cite{rusty} does well, does badly, and does not yet do for this physics
(Section~\ref{sec:ch4-rusty}). A short survey of other domains (Section~\ref{sec:ch4-other}) and exercises close the
chapter.

\begin{warnbox}[: how to read the tiers in this chapter]
Almost nothing in this chapter is Tier~A. The kernel-checked statements of Chapter~2 are about the functions
$g(t)=t^{-s/2}$ (Riesz, $s>0$) and about \emph{admissible} potentials in general. Whether a given laboratory system
\emph{is} described by such a potential is a modelling statement: Tier~C when we argue it, Tier~L when a cited paper
reports it. An observation in a cited experiment is Tier~L; an analogy between two systems is Tier~C; a number we
computed in floating point is Tier~X. A physical system that ``looks triangular'' is never evidence for the
theorem, and a physical system that does not is never evidence against it.
\end{warnbox}

%=====================================================================
\section{Two-dimensional Euler flow and point vortices}\label{sec:ch4-vortices}

\subsection{The vorticity equation}
An ideal (inviscid), incompressible fluid in the plane has velocity $u=(u_1,u_2)$ obeying the Euler equations
\begin{equation}\label{eq:ch4-euler}
\partial_t u + (u\cdot\nabla)u = -\nabla p,\qquad \nabla\cdot u = 0 ,
\end{equation}
with the density set to one. Incompressibility lets us write $u$ through a \emph{stream function}
$\psi$\footnote{The stream function $\psi$ of this section is unrelated to the bond-orientational order parameter
$\psi_6$ used later in this chapter and in Chapter~3.},
\begin{equation}\label{eq:ch4-stream}
u = \nabla^{\perp}\psi := (\partial_y\psi,\,-\partial_x\psi),
\end{equation}
and the scalar \emph{vorticity} $\omega = \partial_x u_2 - \partial_y u_1$ then satisfies
\begin{equation}\label{eq:ch4-poisson}
-\Delta\psi = \omega .
\end{equation}
Taking the curl of \eqref{eq:ch4-euler} eliminates the pressure. In two dimensions the vortex-stretching term
vanishes and one is left with pure transport,
\begin{equation}\label{eq:ch4-vorticity}
\partial_t\omega + u\cdot\nabla\omega = 0 .
\end{equation}
Vorticity is carried by the flow; it is neither created nor destroyed. The whole dynamics is: given $\omega$,
solve the Poisson equation \eqref{eq:ch4-poisson} for $\psi$, form $u$ by \eqref{eq:ch4-stream}, and move
$\omega$ along $u$. This is textbook fluid mechanics \tierC{} (we state it without a citation; any fluid
dynamics text derives it).

\subsection{The Green's function of \texorpdfstring{$-\Delta$}{-Laplacian} in the plane}
To solve \eqref{eq:ch4-poisson} on the whole plane we need a function $G$ with $-\Delta G = \delta_0$.

\begin{proposition}\label{prop:ch4-green}
The function $G(x) = -\frac{1}{2\pi}\log|x|$ satisfies $-\Delta G = \delta_0$ in the sense of distributions on
$\R^2$.
\end{proposition}
\begin{humanproof}
For $x\neq0$ write $r=|x|$. A radial function $f(r)$ has $\Delta f = f'' + f'/r = r^{-1}(r f')'$. With
$f=\log r$, $rf'=1$, so $\Delta \log r = 0$ away from the origin. It remains to find the weight of the delta. Let
$\varphi$ be a smooth test function with compact support. By Green's second identity on the annulus
$\epsilon<|x|<R$ (with $R$ beyond the support of $\varphi$),
\[
\int_{|x|>\epsilon} G\,\Delta\varphi\,dx
= \oint_{|x|=\epsilon}\bigl(G\,\partial_n\varphi - \varphi\,\partial_n G\bigr)\,d\ell ,
\]
where $n$ is the normal pointing \emph{into} the small disk, i.e.\ $\partial_n = -\partial_r$. Now
$\partial_r G = -1/(2\pi r)$, so $-\varphi\,\partial_n G = \varphi\,\partial_r G = -\varphi/(2\pi\epsilon)$ on the
circle, whose length is $2\pi\epsilon$; this term tends to $-\varphi(0)$. The first term is bounded by
$\frac{1}{2\pi}|\log\epsilon|\cdot\|\nabla\varphi\|_\infty\cdot 2\pi\epsilon\to 0$. Hence
$\int G\,\Delta\varphi = -\varphi(0)$, which is $\langle -\Delta G,\varphi\rangle = \varphi(0)$.
\end{humanproof}

The physical reading is Gauss's law: the flux of $-\nabla G$ through any circle around the origin is
$2\pi r\cdot \frac{1}{2\pi r}=1$, independent of $r$. In three dimensions the same computation gives
$1/(4\pi|x|)$; in two it gives a \emph{logarithm}. This single fact is why the planar point-vortex system is the
logarithmic case of our energy problem.

\subsection{Kirchhoff's point vortices}
Concentrate the vorticity at $N$ points: $\omega = \sum_{j=1}^N \Gamma_j\,\delta_{x_j}$, where $\Gamma_j$ is the
\emph{circulation} of vortex $j$ (positive for counter-clockwise rotation). By Proposition~\ref{prop:ch4-green},
\[
\psi(x) = -\frac{1}{2\pi}\sum_{j}\Gamma_j\log|x-x_j| .
\]
Each vortex is moved by the velocity due to all the \emph{others}; its own field is circular around it and does not
move it (this is the standard regularisation, Tier~C). Writing $x_i=(x_i,y_i)$ and
$r_{ij}=|x_i-x_j|$, \eqref{eq:ch4-stream} gives the Kirchhoff--Helmholtz equations
\begin{equation}\label{eq:ch4-kirchhoff}
\dot x_i = -\frac{1}{2\pi}\sum_{j\neq i}\Gamma_j\,\frac{y_i-y_j}{r_{ij}^2},\qquad
\dot y_i = \frac{1}{2\pi}\sum_{j\neq i}\Gamma_j\,\frac{x_i-x_j}{r_{ij}^2}.
\end{equation}

\begin{proposition}[Hamiltonian structure]\label{prop:ch4-ham}
Let
\begin{equation}\label{eq:ch4-H}
H = -\frac{1}{2\pi}\sum_{i<j}\Gamma_i\Gamma_j\log r_{ij}
  = \frac12\sum_{i\neq j}\Gamma_i\Gamma_j\,G(x_i-x_j).
\end{equation}
Then \eqref{eq:ch4-kirchhoff} is equivalent to
$\Gamma_i\dot x_i = \partial H/\partial y_i$ and $\Gamma_i\dot y_i = -\partial H/\partial x_i$ (for $\Gamma_i\ne0$).
\end{proposition}
\begin{humanproof}
$\partial_{y_i}\log r_{ij} = (y_i-y_j)/r_{ij}^2$, and $y_i$ appears in the terms with pair $\{i,j\}$ for every
$j\neq i$. So $\partial H/\partial y_i = -\frac{1}{2\pi}\sum_{j\ne i}\Gamma_i\Gamma_j(y_i-y_j)/r_{ij}^2$; divide by
$\Gamma_i$ and compare with \eqref{eq:ch4-kirchhoff}. The $x$-equation is the same computation with a sign.
\end{humanproof}

So $(x_i,\Gamma_i y_i)$ are canonical coordinates: the plane of each vortex is its own phase space. The
Hamiltonian is conserved. So are the \emph{linear impulse} $P=\sum_i\Gamma_i x_i$, $Q=\sum_i\Gamma_i y_i$, and
the \emph{angular impulse} $I=\sum_i\Gamma_i|x_i|^2$. These three come from translation and rotation invariance;
in fact they are conserved for \emph{any} velocity law of the form
$\dot x_i = \sum_{j\neq i}\Gamma_j\,(x_i-x_j)^{\perp} f(r_{ij})$, a fact the code below uses as a negative control
(Exercise~\ref{ex:ch4-impulse}).

\begin{remark}[Two different divergences]\label{rem:ch4-divergences}
The kinetic energy of the flow is $\frac12\int|u|^2 = \frac12\int\omega\psi$ (integrate by parts). For point
vortices this is
$\frac12\sum_{i,j}\Gamma_i\Gamma_jG(x_i-x_j)$, \emph{including} $i=j$, where $G(0)=+\infty$. Dropping the
diagonal (the \emph{self-energy} of each vortex) gives \eqref{eq:ch4-H}; it is a renormalisation by an infinite
constant that does not depend on positions. A second, independent divergence appears for infinitely many vortices
of one sign at positive density: the pair sum $\sum_{j}\log r_{ij}$ diverges at large distances. The first
divergence is local and harmless for dynamics; the second is the one that puts the problem outside the theorem.
\end{remark}

\subsection{Why this is the case \texorpdfstring{$s=0$}{s = 0}, and why it is outside the theorem}
For like-signed vortices, $\Gamma_i\Gamma_j>0$, the pair energy $-\frac{\Gamma_i\Gamma_j}{2\pi}\log r$
\emph{decreases} with distance: like vortices repel in the energy sense, exactly as like charges do. In the Riesz
family $r^{-s}$, the logarithm is the $s\to 0$ limit after subtracting a constant:
\[
\frac{r^{-s}-1}{s}\;\longrightarrow\;-\log r\qquad (s\to0^+).
\]
In the language of Chapter~1, with $t=r^2$, the logarithmic potential is $g(t)=-\tfrac12\log t$. It fails
\lean{AdmissiblePotential} \tierC{} (an elementary check, Exercise~\ref{ex:ch4-log}): it is negative for $t>1$.
Moreover, every nonzero lattice point has $|a|^2\ge2/\sqrt3>1$, so the real lattice sum
$\sum_{a\ne0}-\tfrac12\log|a|^2$ diverges to $-\infty$. (In Lean, \lean{ENNReal.ofReal} truncates each negative term to
$0$, so \lean{latticeEnergy} of this $g$ is $0$.) Either way, the theorem does not apply (the logarithm is not a
nonnegative completely monotone function). The theorem of Chapter~2 says nothing about it.

What replaces it is a \emph{renormalised} energy, in which the infinite interaction with distant particles is
cancelled by a neutralising background. Upstream's manuscript states, in its Section~8 and after Petrache and
Serfaty~\cite{ps}, that for $0\le s<2$ (the logarithm included) the triangular lattice minimises the renormalised
energy~\cite{oai} \tierL. That argument is in the AI-generated, not peer-reviewed manuscript only; it is \emph{not}
in Lean anywhere. Chapter~5 explains the construction and the open endpoint $s=2$.

\subsection{Rotating frames: the vortex lattice as a logarithmic jellium}
A family of $N$ vortices that rotates rigidly at angular velocity $\Omega$ about the origin satisfies
$\dot x_i = -\Omega y_i$, $\dot y_i=\Omega x_i$. Insert this into Proposition~\ref{prop:ch4-ham}:
$\Gamma_i(-\Omega y_i) = \partial H/\partial y_i$ and $\Gamma_i\Omega x_i = -\partial H/\partial x_i$, that is,
\begin{equation}\label{eq:ch4-rotframe}
\nabla_{x_i}\Bigl(H + \frac{\Omega}{2}\,I\Bigr) = 0\quad\text{for all } i,\qquad
I=\sum_i\Gamma_i|x_i|^2 .
\end{equation}
\emph{Rigidly rotating vortex configurations are exactly the critical points of $H_\Omega=H+\frac{\Omega}{2}I$.}
For $N$ equal vortices $\Gamma>0$ and $\Omega>0$,
\[
H_\Omega = -\frac{\Gamma^2}{2\pi}\sum_{i<j}\log r_{ij} + \frac{\Omega\Gamma}{2}\sum_i|x_i|^2,
\]
a logarithmic repulsion in a harmonic bowl. The bowl is precisely the potential of a uniform neutralising
background. Indeed, let a uniform disk $D$ carry vorticity density $-\Gamma n_b$. A vortex $\Gamma$ at $x\in D$
has interaction energy $V(x) = -\Gamma^2 n_b\int_D G(x-y)\,dy$ with it, and by Proposition~\ref{prop:ch4-green},
$-\Delta_x\int_D G(x-y)\,dy = 1$ inside $D$, so $\Delta V = \Gamma^2 n_b$. On the other hand
$\Delta\bigl(\frac{\Omega\Gamma}{2}|x|^2\bigr) = 2\Omega\Gamma$. Both are radial, so they agree up to a constant
exactly when
\begin{equation}\label{eq:ch4-feynman}
n_b = \frac{2\Omega}{\Gamma}.
\end{equation}
The rotating frame therefore turns the like-sign vortex gas into a \emph{two-dimensional one-component plasma}: a
logarithmic jellium with background density $2\Omega/\Gamma$ \tierC. If the vortex density equals the background
density, the system is neutral on average, and the large-scale divergence of Remark~\ref{rem:ch4-divergences}
cancels. In a quantum superfluid, every vortex carries the same quantised circulation $\Gamma=h/m$, and
\eqref{eq:ch4-feynman} becomes the familiar statement that a rotating superfluid mimics solid-body rotation with a
vortex density $2\Omega m/h$ (this is often called Feynman's rule; we give no citation for it here) \tierC.

\begin{warnbox}[: what the analogy does and does not say]
The identification ``rotating vortices $=$ log-jellium'' is an exact rewriting for \emph{point} vortices in an
\emph{unbounded} ideal fluid. Real vortices have cores, real containers have walls (which add image vortices), and
real superfluids are described by a Gross--Pitaevskii or similar field theory, not by point particles. Moreover,
the renormalised-energy statement for the logarithm is Tier~L through~\cite{oai} only. So ``the triangular vortex
lattice of a rotating superfluid is explained by universal optimality'' is \emph{not} a claim of this book. The
claim is weaker: the energy that a rotating vortex array makes stationary is, after an exact rewriting, the
$s=0$ jellium energy that Chapter~5 discusses \tierC.
\end{warnbox}

\subsection{Experiments: vortex crystals and vortex lattices}

\begin{physbox}[: vortex crystals in a magnetised electron column]
Fine, Cass, Flynn and Driscoll observed the relaxation of two-dimensional turbulence to \emph{vortex crystals}
in a magnetised pure-electron column~\cite{fine} \tierL. Why is an electron column a two-dimensional fluid? In a
strong axial field $B\hat z$, the electrons' guiding centres drift with $v = E\times B/B^2$. With $E=-\nabla\phi$
and $B=B\hat z$, $v = \bigl(-\partial_y\phi,\ \partial_x\phi\bigr)/B$. Compare \eqref{eq:ch4-stream}:
$\psi = -\phi/B$. Poisson's law for the electron density $n$, $\Delta\phi = en/\varepsilon_0$, then gives
$\omega = -\Delta\psi = en/(\varepsilon_0B)$: \emph{the electron density is the vorticity} \tierC{} (the standard
guiding-centre drift argument, not taken from~\cite{fine}). The drift dynamics is the 2D Euler equation
\eqref{eq:ch4-vorticity}, with the conducting wall of the trap fixing the boundary condition on $\phi$.

What~\cite{fine} reports is stated in its title: two-dimensional turbulence relaxes to ``vortex crystals''. We
quote nothing more from that paper here: no number of vortices, no arrangement, no lifetime. The physics we add is
our own reasoning \tierC. Ideal drift dynamics is \emph{conservative}: $H$ is not minimised by dissipation. An
ordered vortex pattern that persists in such a flow is naturally read as a rigidly rotating state, i.e.\ a
critical point of $H_\Omega$ in \eqref{eq:ch4-rotframe}, here in a bounded domain with a wall. The connection
with jellium-type energy minimisation is an analogy \tierC.
\end{physbox}

\begin{physbox}[: vortex lattices in a rotating Bose--Einstein condensate]
Abo-Shaeer, Raman, Vogels and Ketterle observed vortex lattices in rotating Bose--Einstein condensates, the vortices
arranging themselves in a triangular (Abrikosov) pattern~\cite{aboshaeer} \tierL. Unlike the ideal point-vortex
flow, a real condensate has dissipation, so its vortex array can settle into a state of low energy in the rotating
frame \tierC{} (our gloss, not a statement taken from~\cite{aboshaeer}). In the rotating frame, the energy that the
lattice minimises contains the term $-\Omega L_z$, whose effect on the vortices is the harmonic, neutralising bowl
of \eqref{eq:ch4-rotframe}. The resulting ``logarithmic jellium'' picture is an analogy, with all the caveats of the
box above \tierC: cores, the finite and inhomogeneous condensate, and the edge.
\end{physbox}

\subsection{A computational check: conservation in a point-vortex code}\label{sec:ch4-pvcode}
Before any physics is extracted from a simulation, one should check that it conserves what the equations
conserve, and that the check \emph{can fail}. The script below (57 lines,
\texttt{code/ch4\_point\_vortices.py} in the book directory) integrates \eqref{eq:ch4-kirchhoff} with
\texttt{scipy.integrate.solve\_ivp}. It runs three experiments.
\begin{enumerate}
\item \emph{Known answer (positive control).} A regular ring of $n$ equal vortices ($\Gamma=1$) of radius $R$
rotates rigidly at $\Omega = (n-1)\Gamma/(4\pi R^2)$ (Exercise~\ref{ex:ch4-ring}). For $n=5$, $R=1$:
$\Omega = 1/\pi$.
\item \emph{Conservation.} Eight vortices of random signs and strengths, integrated to $t=10$ with an explicit
Runge--Kutta method (DOP853) and with an implicit BDF method; the drifts of $H$, $P$, $Q$, $I$ over 201 output times
are printed.
\item \emph{Negative control.} The same system with the kernel $1/r$ in place of $1/r^2$ in
\eqref{eq:ch4-kirchhoff} (\texttt{power=1}). This is no longer the flow generated by the logarithmic $H$, so $H$
must drift, while $P$, $Q$, $I$ should still be conserved by symmetry.
\end{enumerate}

\begin{lstlisting}[language=Python,caption={\texttt{code/ch4\_point\_vortices.py} (verbatim).}]
"""N point vortices in the plane (Kirchhoff): conservation checks with a known answer and a negative control."""
import numpy as np
from scipy.integrate import solve_ivp
from scipy.optimize import OptimizeResult as OdeResult


def rhs(t: float, z: np.ndarray, gam: np.ndarray, power: int = 2) -> np.ndarray:
    """Velocity of each vortex: u_i = -(1/2pi) sum_j G_j (y_i-y_j)/r^p, v_i = (1/2pi) sum_j G_j (x_i-x_j)/r^p."""
    x, y = z[0::2], z[1::2]
    dx, dy = x[:, None] - x[None, :], y[:, None] - y[None, :]
    r2 = dx**2 + dy**2
    np.fill_diagonal(r2, 1.0)
    k = gam[None, :] / (2 * np.pi * r2 ** (power / 2))
    np.fill_diagonal(k, 0.0)
    out = np.empty_like(z)
    out[0::2], out[1::2] = -(k * dy).sum(1), (k * dx).sum(1)
    return out


def invariants(z: np.ndarray, gam: np.ndarray) -> tuple[float, float, float, float]:
    """Hamiltonian H = -(1/2pi) sum_{i<j} G_i G_j log r_ij, impulses P, Q and angular impulse I."""
    x, y = z[0::2], z[1::2]
    r = np.hypot(x[:, None] - x[None, :], y[:, None] - y[None, :])
    iu = np.triu_indices(len(gam), 1)
    h = -np.sum(np.outer(gam, gam)[iu] * np.log(r[iu])) / (2 * np.pi)
    return h, float(gam @ x), float(gam @ y), float(gam @ (x**2 + y**2))


def drift(sol: OdeResult, gam: np.ndarray) -> list[float]:
    """Max over output times of |invariant(t) - invariant(0)| (relative for H and I)."""
    inv = np.array([invariants(sol.y[:, k], gam) for k in range(sol.y.shape[1])])
    d = np.abs(inv - inv[0]).max(0)
    return [d[0] / abs(inv[0, 0]), d[1], d[2], d[3] / abs(inv[0, 3])]


# Known answer: a ring of n equal vortices (G = 1, radius R) rotates rigidly at omega = (n-1)/(4 pi R^2).
n, R, T = 5, 1.0, 20.0
th = 2 * np.pi * np.arange(n) / n
z0 = np.column_stack([R * np.cos(th), R * np.sin(th)]).ravel()
ring = solve_ivp(rhs, (0, T), z0, args=(np.ones(n),), method="DOP853", rtol=1e-12, atol=1e-12)
angle = np.unwrap(np.arctan2(ring.y[1], ring.y[0]))[-1]
print(f"ring n=5: omega numeric {angle / T:.12f}  exact {(n - 1) / (4 * np.pi * R**2):.12f}")

# Random mixed-sign system, two integrators.
rng = np.random.default_rng(2026)
N = 8
gam = rng.choice([-1.0, 1.0], N) * rng.uniform(0.5, 1.5, N)
z0 = rng.uniform(-1, 1, 2 * N)
t_eval = np.linspace(0, 10, 201)
for method in ("DOP853", "BDF"):
    sol = solve_ivp(rhs, (0, 10), z0, args=(gam,), method=method, rtol=1e-10, atol=1e-12, t_eval=t_eval)
    dH, dP, dQ, dI = drift(sol, gam)
    print(f"{method:6s} nfev={sol.nfev:6d}  rel dH={dH:.2e}  dP={dP:.2e}  dQ={dQ:.2e}  rel dI={dI:.2e}")

# Negative control: a 1/r kernel instead of 1/r^2 is not the gradient of the log Hamiltonian.
bad = solve_ivp(rhs, (0, 10), z0, args=(gam, 1), method="DOP853", rtol=1e-10, atol=1e-12, t_eval=t_eval)
print("control (1/r kernel): rel dH={:.2e}  dP={:.2e}  dQ={:.2e}  rel dI={:.2e}".format(*drift(bad, gam)))
\end{lstlisting}

\begin{numbox}[: output of the point-vortex script (Tier X)]
Command, exactly as run (two cores):
\begin{lstlisting}
taskset -c 2,3 /home/callensxavier_gmail_com/AutoevolveAI/.venv/bin/python -I /home/callensxavier_gmail_com/AutoevolveAI/.claude/worktrees/openai-math-discovery/papers/textbook_universal_optimality/code/ch4_point_vortices.py
\end{lstlisting}
Output, verbatim:
\begin{lstlisting}
ring n=5: omega numeric 0.318309886184  exact 0.318309886184
DOP853 nfev= 57329  rel dH=1.92e-08  dP=7.55e-15  dQ=4.44e-15  rel dI=9.68e-11
BDF    nfev=166250  rel dH=5.78e-07  dP=2.60e-14  dQ=2.02e-14  rel dI=4.05e-09
control (1/r kernel): rel dH=3.61e-02  dP=2.11e-15  dQ=8.88e-16  rel dI=8.21e-11
\end{lstlisting}
\end{numbox}

How to read it \tierX:
\begin{itemize}
\item The ring's rotation rate agrees with $1/\pi$ to the twelve printed digits. The code's velocity law and its
normalisation $1/2\pi$ are therefore right.
\item With DOP853 the linear impulses are conserved to round-off ($\sim10^{-14}$), the angular impulse to about
$10^{-10}$ relative and the Hamiltonian to about $2\times10^{-8}$ relative. $H$ is the most sensitive
invariant here because it contains $\log r_{ij}$, which is large when two vortices of opposite sign come close.
\item BDF conserves everything less well (about $30\times$ worse for $H$) and uses about three times as many
right-hand-side evaluations. This is \emph{not} a verdict on BDF in general. Without a supplied Jacobian, SciPy's BDF
builds one by finite differences, and those evaluations are counted in \texttt{nfev}. More importantly, the
point-vortex flow is Hamiltonian, not dissipative: the eigenvalues of its Jacobian lie on or near the imaginary
axis, the problem is not stiff, and an implicit method built for stiff decay brings cost without benefit. The
gradient flows of Chapter~3 are the opposite case (Section~\ref{sec:ch4-rusty}).
\item The negative control does what it should: with the wrong kernel, $H$ drifts by $3.6\times10^{-2}$, more than
six orders of magnitude above the correct run, while $P$, $Q$ and $I$ remain conserved. A conservation test that
checked only $P$, $Q$, $I$ would have passed the wrong code. Only $H$ tests that the kernel is the gradient of the
logarithm.
\end{itemize}

%=====================================================================
\section{Colloids at fluid interfaces}\label{sec:ch4-colloids}

Colloids are particles of roughly micrometre size suspended in a liquid. Confined to a plane---by gravity onto a
flat glass surface, or by surface tension at a liquid--air interface---and made to repel each other, they form
two-dimensional crystals that can be filmed particle by particle. They are the cleanest laboratory realisation of
the theorem's setting that we know of, for two reasons. The interaction law can be tuned and is close to a pure
power law in one important case; and the dynamics is overdamped motion in a viscous fluid, which is the gradient
flow that our numerics integrate.

\subsection{Superparamagnetic colloids: a Riesz \texorpdfstring{$s=3$}{s = 3} system}
Zahn, Wille, Maret, Sengupta and Nielaba studied superparamagnetic colloids at a flat air--water interface and
measured the elastic properties of the resulting two-dimensional crystals by video
microscopy~\cite{zahn} \tierL. A magnetic field $B$ perpendicular to the interface induces in each particle a
magnetic moment $m\propto\chi B$, where $\chi$ is the particle's susceptibility. All moments are parallel and
perpendicular to the plane, so for two particles at in-plane distance $r$ the dipole--dipole energy is
\begin{equation}\label{eq:ch4-dipole}
U(r) = \frac{\mu_0}{4\pi}\,\frac{m^2}{r^3}
\end{equation}
(the angular factor $1-3\cos^2\theta$ of the dipole law equals one because the moments are perpendicular to their
separation) \tierC. This is a Riesz potential with $s=3$. In the language of Chapter~1,
$g(t) = t^{-3/2}$ up to a positive constant. The Lean results of Chapter~2 are stated for $t^{-s/2}$ itself;
multiplying by a positive constant multiplies every energy by that constant and changes no inequality (our step,
\tierC, and trivial). For $t^{-3/2}$ they say \tierA{} (kernel-checked, conditional on the fidelity of the
definitions):
\begin{itemize}
\item $g$ is admissible (\lean{riesz_admissible}, which needs $s>0$), so the triangular lattice minimises the lower
energy among all locally finite configurations of density one (\lean{triangular_riesz_optimal}, D1);
\item its lattice energy is finite (\lean{latticeEnergy_riesz_lt_top}, which needs $s>2$; D1b);
\item the same holds at any density $\rho$, with the energy scaled accordingly (D2, theorem
\lean{triangular_optimal_any_density}). For $s=3$ the lattice energy at density $\rho$ is $\rho^{3/2}$ times its value at density one
(\lean{riesz_lattice_homogeneous}), and that scaled value is a lower bound for every competitor of density $\rho$
(\lean{riesz_any_density_homogeneous}).
\end{itemize}
The Tier~A content is the statement about $t^{-3/2}$. That the colloids interact by \eqref{eq:ch4-dipole} is a
modelling assumption \tierC, with known corrections: particles have finite size, their moments induce each other
(mutual induction), and the interface may not be perfectly flat.

It is useful to have a single dimensionless number for such a system. At area density $n$ the typical spacing is
$\propto n^{-1/2}$, so the typical dipolar energy is $\propto \mu_0 m^2 n^{3/2}$, and the ratio to the thermal
energy is a coupling of the form
\[
\Gamma_{\rm mag}\;\propto\;\frac{\mu_0\,(\chi B)^2\, n^{3/2}}{k_BT}.
\]
Raising $B$ at fixed $n$ cools the system in the only sense that matters: it increases the coupling. A slow
ramp of $B$ is the laboratory analogue of the simulated annealing of Chapter~3 \tierC.

\subsection{Charged colloids in cavities}
Bubeck, Neser, Bechinger and Leiderer studied small two-dimensional colloidal crystals confined in circular,
hexagonal and square cavities~\cite{bubeck} \tierL. The tiered paper~\cite{tiered}, summarising the publisher's
abstract (its full text was read neither there nor here), reports that these clusters contained about 10 to 100
particles and that their structure was dominated by the shape of the confining region. For us, such cavities are the experimental counterpart of the ``closed disk'' of
Chapter~3, in which a soft wall holds the particles. Edge effects are not a nuisance here; they are the phenomenon.
A cavity of $N$ particles has a boundary layer whose share of particles decreases only like $N^{-1/2}$.

\subsection{Hydrodynamics: from Stokes drag to the gradient flow}\label{sec:ch4-stokes}
Why does a colloidal monolayer relax by the gradient flow $\dot x = -\nabla E$ used in Chapter~3? Because it moves
through a viscous liquid at very small Reynolds number.

\paragraph{Stokes drag.} A sphere of radius $a$ moving at velocity $v$ through an unbounded liquid of viscosity
$\eta$, slowly enough that inertia of the liquid is negligible, feels the drag $F_{\rm drag}=-\gamma v$ with
$\gamma = 6\pi\eta a$ (Stokes' law, textbook result, no-slip boundary condition) \tierC. A particle sitting
\emph{at} an air--water interface is only partly immersed, and its drag coefficient differs from the bulk value;
we leave $\gamma$ symbolic and quote no number.

\paragraph{The Langevin equation.} Newton's law for particle $i$, of mass $m$, at temperature $T$ reads
\begin{equation}\label{eq:ch4-langevin}
m\,\ddot x_i = -\gamma\,\dot x_i - \nabla_{x_i}E(x_1,\dots,x_N) + \xi_i(t),\qquad
\langle\xi_i(t)\otimes\xi_j(t')\rangle = 2\gamma k_BT\,\delta_{ij}\,\mathrm{Id}\;\delta(t-t'),
\end{equation}
where $\xi$ is the random force of the solvent molecules. Its strength is fixed by the fluctuation--dissipation
theorem: with this choice, the stationary distribution is the Boltzmann distribution $\propto e^{-E/k_BT}$ \tierC.

\paragraph{The overdamped limit.} Without the forces $-\nabla E$ and $\xi$, the velocity decays as
$e^{-t/\tau_m}$ with $\tau_m = m/\gamma$. For a micrometre colloid in water, $\tau_m$ is many orders of magnitude
shorter than the time the particle needs to move by a fraction of its spacing (Exercise~\ref{ex:ch4-taum}, where
you estimate it yourself). On time scales $\gg\tau_m$ the inertial term is negligible, and
\eqref{eq:ch4-langevin} becomes the \emph{overdamped Langevin equation}
\begin{equation}\label{eq:ch4-overdamped}
\gamma\,\dot x_i = -\nabla_{x_i}E + \xi_i ,\qquad D=\frac{k_BT}{\gamma}\ \ \text{(Einstein relation).}
\end{equation}
At $T=0$ it is the gradient flow $\dot x = -\gamma^{-1}\nabla E$. Rescaling time by $\gamma$ gives exactly the
equation $\dot x = -\nabla E$ that rusty-SUNDIALS integrates with CVODE in Chapter~3. The annealing of Chapter~3 is
modelled on \eqref{eq:ch4-overdamped} at slowly decreasing $T$, but the code uses a capped Euler--Maruyama scheme
(heuristic annealing, see Chapter~3), which is not a consistent discretisation of it \tierC.

\paragraph{Hydrodynamic interactions.} The model above treats each particle as feeling only its own drag.
In reality, a moving colloid sets the liquid around it in motion, and that flow drags its neighbours. To leading
order this is described by a \emph{mobility tensor} $M_{ij}$ (Oseen or Rotne--Prager forms; see the literature),
and the zero-temperature dynamics becomes $\dot x_i = -\sum_j M_{ij}\nabla_{x_j}E$. Two facts survive \tierC:
\begin{enumerate}
\item If $M$ is symmetric positive definite, $\frac{d}{dt}E = -\sum_{ij}\nabla_iE\cdot M_{ij}\nabla_jE\le0$: the
energy still decreases.
\item The stationary points ($\nabla E = 0$) and their stability (the Hessian) do not depend on $M$.
\end{enumerate}
Hydrodynamics therefore changes the \emph{paths} to the minimum and the efficiency of annealing, not the
minimum. This is why the question ``what is the ground state?''\ belongs to the energy alone. It is also why the
theorem, which is about the energy alone, is relevant to colloids at all.

%=====================================================================
\section{Plasmas}\label{sec:ch4-plasmas}

\subsection{Dusty-plasma crystals: the Yukawa potential}
Thomas, Morfill, Demmel, Goree, Feuerbacher and M\"ohlmann observed Coulomb crystallisation of charged
microparticles levitated in a plasma~\cite{thomas} \tierL. The grains acquire a large negative charge from the
plasma, and the surrounding ions and electrons screen it. In the simplest (Debye--H\"uckel) model the screened pair
potential is of Yukawa form,
\[
V(r) = \frac{e^{-\kappa r}}{r}\quad\text{(up to a positive constant; $\kappa$ the inverse screening length)}
\]
\tierC. As a function of $t=r^2$, $g(t) = e^{-\kappa\sqrt t}/\sqrt t$. Is it admissible? The function $1/\sqrt t$
is completely monotone; $e^{-\kappa\sqrt t}$ is completely monotone because $e^{-\kappa u}$ is completely monotone
in $u$ and $\sqrt t$ has a completely monotone derivative; products of completely monotone functions are
completely monotone. So $g$ satisfies the hypotheses of the upstream theorem \tierC{} (Exercise~\ref{ex:ch4-yukawa}).
This step is \emph{not} in our Lean files, which is why every statement about Yukawa crystals in this book is
Tier~C at best. The numerical study of Chapter~3 is for exactly this potential (with a shifted-force cutoff at
$r_c=8$): at $\kappa=2$, with $N=288$ in a periodic box at density one, the reference energies per particle are
$e_{\rm tri}=0.3935333668$ and $e_{\rm sq}=0.4009647699$ (each pair counted once, as in Chapter~3: one half of
Chapter~1's $E_g$), and none of 100 relaxations (20 of them
incommensurate-box controls) fell below $e_{\rm tri}$ \tierX.

Real dusty plasmas leave the theorem's class in several ways \tierC{} (general knowledge, not cited here). The ion
flow past a grain builds a wake that attracts other grains, chiefly those below it; the wake force is non-reciprocal,
so there is no pair energy at all. Grains move out of the plane; and far from a grain, the potential is not Yukawa. These confounds are listed again in
protocol L2 below.

\subsection{Electrons on liquid helium}
Grimes and Adams observed a liquid-to-crystal phase transition of a classical two-dimensional sheet of electrons held
above the surface of liquid helium~\cite{grimes} \tierL. The pair interaction is Coulomb,
$e^2/(4\pi\varepsilon_0 r)$. Here `Coulomb' means the three-dimensional $1/r$ law restricted to the plane (Riesz
$s=1$); the planar Coulomb (logarithmic, $s=0$) interaction is distinguished where it matters. It is Riesz with $s=1$, so the infinite lattice sum diverges and the system is neutralised
by positive charge on electrodes \emph{below} the sheet. This is jellium in the sense of Chapter~5---but with a
caveat that the clean mathematical jellium does not have.

\paragraph{Image charges change the exponent.} Model the electrode as a grounded conducting plane a distance $d$ below
the sheet. Each electron then has an image of opposite charge ($+e$ for an electron) at distance $2d$ from it, vertically below. Two electrons
at in-plane distance $r$ interact directly and through each other's images:
\[
U(r) = \frac{e^2}{4\pi\varepsilon_0}\Bigl(\frac1r - \frac{1}{\sqrt{r^2+4d^2}}\Bigr)
\;\approx\;\frac{e^2}{4\pi\varepsilon_0}\,\frac{2d^2}{r^3}\qquad (r\gg d)
\]
\tierC{} (Exercise~\ref{ex:ch4-image}). At distances large compared with $d$, the interaction becomes
\emph{dipolar}, $s=3$, and the system leaves the jellium regime and enters the finite-energy class of the theorem.
Electrons on helium are therefore jellium only when $d$ is much larger than the interparticle spacing.

\subsection{Wigner crystal energetics}
For the clean 2D Wigner crystal (point charges with $1/r$ interaction in a plane, uniform neutralising background
in the same plane) the energies of the triangular and square lattices were computed by Bonsall and
Maradudin~\cite{bm}, who found the triangular lattice lower and the square lattice dynamically unstable (an
imaginary phonon branch) \tierL. Alves, Bendazzoli, Evangelisti and Berger gave ten-digit Madelung energies by a
real-space method~\cite{alves} \tierL.

\begin{numbox}[: jellium controls from Chapter~3 (Tier X)]
Our Ewald implementation (rusty-SUNDIALS branch \lean{physics/jellium-wigner-2026-10-11}) reproduces the published
values: triangular energy per particle $-1.960515789320$ against the literature value $-1.960515789825$ (relative
difference $2.6\times10^{-10}$); square $-1.950132460$, relative difference $4.3\times10^{-10}$ from the
literature; the energy is invariant to better than $10^{-10}$ over a $4\times$ range of the Ewald splitting
parameter; and the square lattice is a saddle point with 64 unstable modes, consistent with~\cite{bm}. The main
relaxation study: \JelliumMainStatus.
\end{numbox}

These are controls, not results: they show that the code computes the energy that the literature computes. They
say nothing about whether some non-lattice configuration could do better; that is the question the theorem
answers (for $s>0$ in the finite-energy class) and that Chapter~5 discusses for $s=1$.

%=====================================================================
\section{Laboratory protocols L1--L4}\label{sec:ch4-lab}

The tiered paper~\cite{tiered} proposes four experiments. We summarise them here with their caveats. They are
\emph{proposals} \tierC; none has been run. An experiment cannot test the Lean statement, which is about exact
potentials and infinite configurations. What it can test is whether real two-dimensional repulsive systems behave as
the theorem and the numerics suggest.

All four share one analysis script, \texttt{lab/lab\_analysis.py} (in the tiered paper's directory), which takes
tracked particle coordinates and computes the same observables as the simulations:
\begin{itemize}
\item the \emph{density}, from the mean Voronoi cell area of the same interior particles over which the energy is
averaged;
\item the \emph{coordination}, by the $1.35\,d_0$ neighbour rule, and the local and global bond-orientational order
$|\psi_6|$;
\item the \emph{energy} per interior particle for a stated interaction law, truncated at $r_c$ with a uniform tail
correction, compared with a triangular lattice at the same density and the same truncation.
\end{itemize}
The script's self-test passes: a perfect lattice gives a gap at round-off level, thermal noise raises the energy
above the lattice, and a square lattice is detected and placed above it. Two caveats bear on every kill criterion.
The truncated potential is not exactly completely monotone, so the comparison is between equally truncated
energies. And because at $s=3$ the energy scales like $\rho^{3/2}$, a relative density error $\delta\rho/\rho$
becomes a relative energy error of about $1.5\,\delta\rho/\rho$ (Exercise~\ref{ex:ch4-density}), which must enter
the error bar. Each protocol should be preregistered with the script, a predicted bracket and its kill criterion.

\paragraph{L1. Superparamagnetic colloids at a flat air--water interface (Riesz $s=3$).} Following~\cite{zahn}.
\emph{Procedure:} ramp $B$ slowly from the fluid to the crystal at fixed area fraction; track positions by video
microscopy; record $\chi$ and $B$ so that the prefactor of \eqref{eq:ch4-dipole} is known. \emph{Predictions:}
(a) at large coupling, $|\psi_6|\to1$ and the interior defect fraction $\to0$; (b) the measured dipolar energy per
interior particle exceeds the triangular value at the measured density, by an amount that shrinks as $B$ grows
(the thermal excess); (c) a particle removed with optical tweezers leaves a vacancy that relaxes locally, without
global restructuring; (d) a square arrangement imposed with holographic tweezers and released relaxes away from
square---the analogue of the square-saddle control computed for Yukawa and jellium (the $s=3$ square lattice itself
was not computed). \emph{Kill criterion:} a slowly ramped, defect-free interior whose dipolar energy lies below the
triangular value beyond the measurement error, with the law's prefactor independently confirmed.

\paragraph{L2. Dusty-plasma monolayer (Yukawa).} Following~\cite{thomas}. The radial confinement is close to
parabolic, the analogue of the ``open bowl'' of Chapter~3. \emph{Procedure:} measure $\kappa$ and the coupling from
the phonon spectrum or the pair correlation; image the monolayer; run the shared script; simulate the bowl model at
the measured $\kappa$ and $N$. \emph{Prediction:} a near-triangular interior with defects concentrated at the rim;
the interior defect fraction is compared with simulation at matched $\kappa$ and $N$. \emph{Confounds} \tierC{} (general knowledge, not cited here): ion-wake
attraction (non-reciprocal, chiefly between vertically adjacent grains, able to drive a mode-coupling instability
in a monolayer), out-of-plane motion, and departures from a pure Yukawa law at large distance all move the system
outside the theorem's class. Neutral-gas drag damps the dynamics but does not change the interaction law.

\paragraph{L3. Charged colloids in circular, hexagonal and square cavities (closed disk).} Following~\cite{bubeck}.
\emph{Procedure:} repeat the circular-cavity case at $N\approx300$, image the interior, compare with closed-wall
simulations at matched $N$ and wall stiffness. \emph{Prediction:} the edge defect fraction exceeds the interior one,
as in the closed-disk simulations of Chapter~3. Compared with L2 at matched $N$ and density, this is the
experimental open-versus-closed comparison---with the confounds that the simulated comparison already has
(non-uniform density in the bowl, compression against the wall, different edge-layer sizes).

\paragraph{L4. Electrons on liquid helium (jellium, $s=1$).} Following~\cite{grimes}. The neutralising charge sits
on electrodes a distance $d$ below the sheet, so this is the jellium case only when $d$ is much larger than the
spacing; otherwise, as computed above, image charges make the interaction dipolar at distances $\gg d$.
\emph{Prediction:} the crystal is triangular. The evidence in~\cite{grimes} is spectroscopic, not a direct image \tierC{} (our reading of~\cite{grimes}), so
a test compares spectroscopic signatures with the triangular Madelung and phonon structure~\cite{bm}. Direct
imaging would need a detection method that the protocol does not specify.

\begin{warnbox}[: what a positive outcome would mean]
Suppose L1 is carried out and every prediction holds. That would show that a real dipolar monolayer behaves as the
ideal model predicts. It would \emph{not} confirm the theorem (which needs no confirmation of that kind and admits
none), and it would \emph{not} confirm the model (a defect-free triangular interior is also what many non-ideal
interactions produce). The kill criterion is asymmetric on purpose: only an undercut of the triangular energy,
with the law independently confirmed, would be informative, and it would point at the model, not at the proof.
\end{warnbox}

%=====================================================================
\section{rusty-SUNDIALS: implications and opportunities}\label{sec:ch4-rusty}

rusty-SUNDIALS~\cite{rusty} is a Rust re-implementation of parts of the SUNDIALS suite. The numerical study of
Chapter~3 (crate \texttt{qf-yukawa-crystal}) uses its CVODE. This section asks what the physics of this chapter
demands from such a library, and what the library currently provides. Every statement about the code below was
checked against the files named; the proposals are \tierC.

\subsection{Stiffness of particle relaxation}
Linearise the gradient flow $\dot x = F(x) = -\nabla E(x)$ around a minimum $x^\ast$:
$\dot{\delta x} = -\mathcal H\,\delta x$, with $\mathcal H$ the Hessian. If $\lambda_1\le\dots\le\lambda_{2N}$ are
its eigenvalues, the Jacobian of the ODE has eigenvalues $-\lambda_k$. In a crystal:
\begin{itemize}
\item the two translation modes have $\lambda=0$ (in our runs, $|\lambda|\le1.5\times10^{-15}$ at $\kappa=2$)
\tierX;
\item the long-wavelength phonons have small $\lambda$; for a crystal their smallest value decreases as the box
grows, because the softest wavenumber scales like $1/L$ \tierC. At the single size studied ($N=288$), the perfect
lattice's smallest non-zero eigenvalue was $0.0236$ \tierX. Separately, over the 80 non-control relaxed runs at
$\kappa=2$, most of them defected, every non-translation eigenvalue was at least $5.4\times10^{-5}$ \tierX; such a
low floor reflects soft modes of defected configurations, not a size effect;
\item the short-wavelength, nearest-neighbour modes have the largest eigenvalue $\lambda_{\max}$.
\end{itemize}
An explicit method is stable only for step sizes $h\lesssim 2/\lambda_{\max}$, while relaxation along the softest
mode needs a time $\sim 1/\lambda_{\min}$. The number of explicit steps therefore grows like the
\emph{stiffness ratio} $\lambda_{\max}/\lambda_{\min}$, which in a crystal grows with the system size (the softest
phonon wavenumber scales like $1/L$) \tierC. We give no numerical value for the ratio because $\lambda_{\max}$ is not
among the recorded results. BDF methods, whose stability regions contain most or all of the left half-plane, take
steps limited by accuracy instead of stability, and this is why CVODE's BDF is the right tool for the quench.

Contrast the point vortices of Section~\ref{sec:ch4-pvcode}: there the dynamics is Hamiltonian, the eigenvalues
are near the imaginary axis, and explicit (or, better, symplectic) integrators are the right tool. The same library
must serve both regimes.

\subsection{BDF versus Adams in rusty-SUNDIALS}
CVODE offers two families: BDF for stiff problems, Adams--Moulton for non-stiff ones. In rusty-SUNDIALS the
Adams path does not work as an Adams method. The crate description of \texttt{qf-bao-distances} records it in one
sentence:
\begin{lstlisting}
... integrated with this repository's own CVODE (BDF; its Adams path was measured to be stuck at order 1) ...
\end{lstlisting}
(\texttt{crates/qf-bao-distances/Cargo.toml}). The README of \texttt{crates/sundials-mcp} gives the measured detail:
\texttt{Method::Adams} never leaves order 1 because a coefficient routine (\texttt{compute\_l}) is a placeholder,
and its error on a smooth quadrature problem scales like $\sqrt{\mathrm{rtol}}$ (it reports $2.1\times10^{-4}$
relative at $\mathrm{rtol}=10^{-7}$). The same README records that many output times from one BDF run degrade
accuracy compared with a fresh solve per output ($4\times10^{-6}$ against $6\times10^{-7}$). These are that
README's measurements, not ours \tierX.

\emph{Implications.} The relaxations of Chapter~3 are stiff and use BDF, so they are not affected by the Adams
defect. A point-vortex or any other Hamiltonian problem would want Adams (or a dedicated symplectic method), and
today it would silently get a first-order method if it asked for Adams. \emph{Opportunity:} implement the
variable-order Adams coefficients and test them against a problem with a known high-order convergence rate
(Project~P2 below).

\subsection{Analytic versus finite-difference Jacobians}
A BDF step solves a nonlinear system by Newton's method and needs the Jacobian $\partial F/\partial x$. For the
gradient flow this is $-\mathcal H$. The two crates of Chapter~3 make different choices, as their source says:
\begin{lstlisting}
crates/qf-yukawa-crystal/src/lib.rs:
/// Overdamped relaxation `dx/dt = F(x)` with CVODE (BDF, analytic Jacobian) until `max |F| < ftol` or `t > tmax`.
crates/qf-yukawa-crystal/src/jellium.rs:
    /// Overdamped relaxation with CVODE (BDF, finite-difference Jacobian) until `max |F| < ftol` or `t > tmax`.
\end{lstlisting}
For the short-range Yukawa potential with a cutoff, the Hessian is a sparse sum of $2\times2$ blocks and is cheap to
write down; the Yukawa relaxation passes it to CVODE (\texttt{.jacobian(jac)} with entries $-\mathcal H$). For
jellium, the Ewald sum makes an analytic Hessian laborious, and CVODE forms the Jacobian by finite differences of
the force: the difference-quotient loop in \texttt{crates/cvode/src/solver.rs} makes one extra right-hand-side
call per column, $J_{:,j}=(f(y+\varepsilon e_j)-f(y))/\varepsilon$, i.e.\ $2N$ Ewald force evaluations per
Jacobian for $N$ particles. The Hessian used for
the stability analysis of jellium is likewise ``central finite differences of the analytic force (symmetrised)''
(source comment in \texttt{jellium.rs}).

\emph{Opportunities} \tierC: an analytic Ewald Hessian (the real-space and reciprocal-space parts can both be
differentiated twice in closed form); automatic differentiation; or a Jacobian-free Newton--Krylov solver in which
only Hessian--vector products are needed, which a single extra force evaluation approximates.

\subsection{Constraints and DAEs: what IDA would add}
Several physical set-ups of this chapter are \emph{constrained}: particles confined by a hard wall (rather than the
soft wall $\varepsilon(r-R)^4$ of Chapter~3), particles on a curved interface, vortices in a bounded domain with
explicit boundary conditions, or relaxation at a fixed density or fixed angular impulse. The natural formulation is
a differential-algebraic system $F(t,y,\dot y)=0$, which is what SUNDIALS' IDA solves. In rusty-SUNDIALS,
\texttt{crates/ida} is not that solver yet; its own header says so:
\begin{lstlisting}
/// A fixed-step Backward Euler prototype for DAE residuals, NOT a general IDA solver.
///
/// It does one fixed-point iteration per step (no Newton-Krylov, no Nordsieck array, no variable
/// order/step size, no error control beyond the fixed-point residual). It approximates `dF/dy ~ 1`
/// (see the fixed-point loop in `solve`), which is only valid for mildly stiff residuals. Treat it
/// as a placeholder for a real IDA port, not as a replacement for it; `CvodeSolver` in the `cvode`
/// crate is the tested, variable-order solver in this repository.
\end{lstlisting}
(\texttt{crates/ida/src/solver.rs}, lines 3--9). It was not used in Chapter~3. \emph{Opportunity:} a real IDA port
(variable order and step, Newton iteration with the residual Jacobian, error control), validated on index-1
problems with closed-form solutions before any physics is run on it.

\subsection{KINSOL: solving \texorpdfstring{$F(x)=0$}{F(x) = 0} directly}
The gradient flow finds \emph{minima}: almost every initial condition flows away from saddle points. But some of
the most informative configurations are saddles. The square lattice is one (Chapter~3 found it a saddle with 72
negative Hessian eigenvalues at $\kappa=2$ and 64 at $\kappa=1$) \tierX; transition states between defect
configurations are others; and rigidly rotating vortex configurations, the critical points of $H_\Omega$ in
\eqref{eq:ch4-rotframe}, are generally saddles too. All of these are zeros of a force, $F(x)=0$. SUNDIALS solves
such systems with KINSOL (Newton and Newton--Krylov with globalisation). rusty-SUNDIALS has no KINSOL: no crate of
that name exists, and a search of the crates' sources and manifests for ``kinsol'' finds nothing. \emph{Opportunity}
\tierC: a KINSOL port, tested on the square-lattice saddle (whose existence and index Chapter~3 already measured)
and on the regular vortex ring of Section~\ref{sec:ch4-pvcode} in its rotating frame.

\subsection{The sundials-mcp server: solvers for AI agents}
\texttt{crates/sundials-mcp} is a Model Context Protocol server (JSON-RPC over standard input and output) through
which AI agents can run rusty-SUNDIALS solvers. Its README lists the tools \texttt{about},
\texttt{list\_problems}, \texttt{solve} (named CVODE problems: \texttt{exponential}, \texttt{robertson},
\texttt{vanderpol}, \texttt{lorenz}, \texttt{domain\_exit}), \texttt{pgpe\_run}, \texttt{cmb\_bound} and
\texttt{bao\_distances}. Two of its design rules are worth copying in any scientific tool that an agent will call:
\begin{itemize}
\item every result carries a field \texttt{ran}, and \texttt{ran: false} means that nothing was computed---never a
pass;
\item a solver failure is reported as a tool error with the solver's own message and \emph{no} partial trajectory;
each known-answer check names an external reference.
\end{itemize}
The server also runs every solve in a worker subprocess, because the CVODE code prints diagnostics to standard
output on some failure paths, which would corrupt the protocol stream. \texttt{pgpe\_run} runs a projected
Gross--Pitaevskii field model, in which quantised vortices can occur, but no tool simulates point particles or
point vortices, and we did not check whether either GPE crate (\texttt{qf-pgpe}, \texttt{qf-gpe2d}) supports a
rotating frame. \emph{Opportunity} \tierC: a \texttt{relax\_crystal} tool (Yukawa or jellium relaxation with the
Chapter~3 controls built in) and a \texttt{point\_vortices} tool (with the ring known-answer test and the
conservation report of Section~\ref{sec:ch4-pvcode}), each returning its controls alongside its result.

\subsection{Student projects}
\begin{enumerate}[label=\textbf{P\arabic*.}]
\item \textbf{Point vortices with CVODE.} Port \texttt{rhs} of Section~\ref{sec:ch4-pvcode} to Rust and integrate it
with rusty-SUNDIALS' CVODE (BDF, the only reliable method today). Reproduce the ring rotation rate and the
conservation table. Compare with the SciPy output above; explain any difference in the drift of $H$. Add the
$1/r$-kernel control. (No CVODE run of this problem exists yet; this is a proposal.)
\item \textbf{Fix Adams.} Implement the variable-order Adams--Moulton coefficients in the CVODE crate and show
order $>1$ on a smooth problem (e.g.\ $y'=-y$), by plotting error against tolerance. Then rerun P1 with Adams.
\item \textbf{Vortex crystals in a rotating frame.} Minimise $H_\Omega$ of \eqref{eq:ch4-rotframe} for $N$ equal
vortices by gradient flow with CVODE (this is a stiff relaxation, like Chapter~3), measure $|\psi_6|$ in the
interior, and check that the interior density approaches $2\Omega/\Gamma$ of \eqref{eq:ch4-feynman}.
\item \textbf{A Newton solver for saddles.} Write a small Newton iteration (a ``KINSOL-lite'') for $F(x)=0$ and
converge onto the square-lattice saddle of the Yukawa crystal from a slightly perturbed square start; count the
negative Hessian eigenvalues and compare with Chapter~3.
\item \textbf{An analytic Ewald Hessian.} Derive and code it, test it against the existing finite-difference
Hessian, and compare the CVODE cost of the jellium relaxation with both Jacobians.
\item \textbf{Hydrodynamic interactions.} Add a Rotne--Prager mobility (see the literature for its form) to the
Yukawa gradient flow and verify numerically that the relaxed configurations and their energies do not change while
the relaxation paths do.
\item \textbf{An MCP tool.} Add a \texttt{point\_vortices} tool to \texttt{sundials-mcp} that follows its honesty
conventions, with the ring as the known-answer check.
\end{enumerate}

%=====================================================================
\section{Other domains, briefly}\label{sec:ch4-other}

\paragraph{Ion crystals.} Ions held in Penning traps can form single-plane crystals. The ions interact by the
three-dimensional Coulomb law (Riesz $s=1$ in the plane), and the trap together with the plasma's rotation provides a
harmonic radial confinement, which, as in \eqref{eq:ch4-rotframe}, plays the role of a neutralising background.
These planar ion crystals are the $s=1$ analogue of the rotating vortex lattice; triangular order is expected for
them \tierC{} (general knowledge; we cite no source here).

\paragraph{Abrikosov lattices in type-II superconductors.} In a type-II superconductor in a magnetic field between
the two critical fields, flux enters as quantised vortices, which arrange themselves in a lattice. Near the upper
critical field, Ginzburg--Landau theory reduces the choice of lattice to minimising a single number (the Abrikosov
parameter), and the triangular lattice is the minimiser; the square lattice that was first proposed is higher
\tierC{} (standard physics, no citation here). The mathematics behind such statements is close to this book's:
Gaussian lattice sums. Montgomery proved that among lattices of fixed covolume the triangular lattice minimises the
theta function $\sum_{a}e^{-\pi\alpha|a|^2}$ for every $\alpha>0$~\cite{montgomery} \tierL, and Rankin and Cassels
studied the analogous problem for the Epstein zeta function~\cite{rankin,cassels} \tierL. These are \emph{lattice}
results: they compare lattices with lattices. The theorem of Chapter~2 compares the triangular lattice with
\emph{every} configuration of density one, periodic or not. Which lattice sum, exactly, expresses the Abrikosov
parameter, and how the result relates to Montgomery's theorem, we leave to the literature \tierC.

\paragraph{Foams and honeycombs are a different problem.} Soap froths in a plane and bees' honeycombs show hexagonal
cells, and it is tempting to file them with the triangular lattice, whose Voronoi cells are regular hexagons. But the
underlying optimisation is different: one minimises the total \emph{perimeter} of a partition of the plane into
cells of equal area. There are no point particles and no pair potential. The ``honeycomb conjecture'' that the
regular hexagonal partition is optimal was proved by Hales \tierC{} (general knowledge; no citation here). Neither that theorem nor the
universal optimality theorem implies the other; they share an answer, not a problem.

%=====================================================================
\section*{Summary}
\begin{itemize}
\item Planar point vortices interact through the 2D Green's function, i.e.\ logarithmically: the case $s=0$, which
lies outside the finite-energy class and needs renormalisation (Chapter~5). In a rotating frame the vortex energy
becomes a logarithmic jellium with background density $2\Omega/\Gamma$ \tierC.
\item Vortex crystals in electron columns~\cite{fine} and vortex lattices in rotating condensates~\cite{aboshaeer}
are observations \tierL; their connection with energy minimisation is an analogy \tierC.
\item Superparamagnetic colloids~\cite{zahn} realise a Riesz $s=3$ law, inside the theorem's class \tierA{} for
the potential, \tierC{} for the modelling. Their overdamped motion is the gradient flow of Chapter~3.
\item Dusty plasmas~\cite{thomas} are Yukawa (Tier~C admissibility); electrons on helium~\cite{grimes} are jellium
only when the electrodes are far, and dipolar otherwise.
\item rusty-SUNDIALS handles stiff relaxation with BDF well; its Adams path is stuck at order~1, its IDA is a
prototype, it has no KINSOL, and its MCP server has no particle tools. Each gap is a student project.
\end{itemize}

%=====================================================================
\section*{Exercises}

\begin{exercise}\label{ex:ch4-green}
Repeat the proof of Proposition~\ref{prop:ch4-green} in $\R^3$ and show that $G(x)=1/(4\pi|x|)$. In $\R^d$,
$d\ge3$, show that $G(x)=c_d|x|^{2-d}$, and find $c_d$ in terms of the surface area of the unit sphere. Which Riesz
exponent $s$ is the Green-function (``Coulomb'') exponent in dimension $d$? Conclude that the \emph{planar} Coulomb
interaction, the Green function of the plane, is the case $s=0$ (logarithm), whereas the $1/r$ law that this book
calls `Coulomb' (the three-dimensional law restricted to the plane) is Riesz $s=1$.
\end{exercise}

\begin{exercise}\label{ex:ch4-ring}
Show that $n$ equal vortices of circulation $\Gamma$ at the vertices of a regular $n$-gon of radius $R$ rotate
rigidly at $\Omega=(n-1)\Gamma/(4\pi R^2)$. (Hint: by symmetry compute the velocity of one vortex; use
$\sum_{k=1}^{n-1}\frac{1}{1-e^{2\pi i k/n}}=\frac{n-1}{2}$.) Check the case $n=2$ directly.
\end{exercise}

\begin{exercise}\label{ex:ch4-impulse}
Let $\dot x_i = \sum_{j\neq i}\Gamma_j\,(x_i-x_j)^{\perp}f(r_{ij})$ with $(a,b)^\perp=(-b,a)$ and any function $f$.
Prove that $P$, $Q$ and $I$ are conserved. Show that for $N=2$ the distance $r_{12}$, hence the logarithmic $H$,
is conserved for \emph{every} $f$. Then explain why, for $N=8$, the negative-control line of the point-vortex output
shows a drift of $H$ but not of $P$, $Q$, $I$. What would a negative control with $N=2$ have shown?
\end{exercise}

\begin{exercise}\label{ex:ch4-log}
Let $g(t)=-\frac12\log t$. Show that $(-1)^r g^{(r)}(t)\ge0$ for all $r\ge1$ and $t>0$, but that $g(t)<0$ for
$t>1$. Which clause of \lean{AdmissiblePotential} fails? Show that adding a constant cannot repair it. Show also
that $\sum_{a\in\Lat\setminus0}g(|a|^2)$ diverges to $-\infty$ (all $|a|^2>1$), and that the Lean \lean{latticeEnergy}, which applies \lean{ENNReal.ofReal} termwise, is $0$.
\end{exercise}

\begin{exercise}\label{ex:ch4-rot}
Derive \eqref{eq:ch4-rotframe}. Then show that for equal vortices the confining term $\frac{\Omega\Gamma}{2}|x|^2$
equals, up to a constant, the interaction energy of a vortex with a uniform disk of vorticity $-\Gamma n_b$
(containing it), with $n_b=2\Omega/\Gamma$. Use Proposition~\ref{prop:ch4-green} and the radial form of the
Laplacian.
\end{exercise}

\begin{exercise}\label{ex:ch4-drift}
Show that the $E\times B$ drift velocity in a uniform field $B\hat z$ is divergence-free when $\phi$ is smooth, and
that the electron density is advected by it. Write the analogue of \eqref{eq:ch4-vorticity} for the density.
Which boundary condition does a grounded cylindrical wall impose, and what does it do to a single off-centre
vortex (method of images)?
\end{exercise}

\begin{exercise}\label{ex:ch4-taum}
Using $\gamma=6\pi\eta a$, estimate the momentum relaxation time $\tau_m=m/\gamma$ of a silica sphere of radius
$1\,\mu$m in water, and the time it takes to diffuse a distance equal to its own radius ($a^2/D$, $D=k_BT/\gamma$)
at room temperature. Look up the material constants yourself and state your sources. Justify the overdamped limit.
\end{exercise}

\begin{exercise}\label{ex:ch4-lyap}
Show that along $\dot x = -M\nabla E$, with $M$ symmetric positive definite, $E$ is non-increasing, and that
the linear stability of a critical point (its number of unstable directions) is the number of negative
eigenvalues of the Hessian, independently of $M$. (Hint: $M^{1/2}\mathcal H M^{1/2}$ and $\mathcal H$ have the same
inertia.)
\end{exercise}

\begin{exercise}\label{ex:ch4-yukawa}
Prove that $g(t)=e^{-\kappa\sqrt t}/\sqrt t$ is completely monotone on $(0,\infty)$ for $\kappa\ge0$. Use: (i) if
$\varphi$ is completely monotone and $h$ is positive with a completely monotone derivative, then
$\varphi\circ h$ is completely monotone; (ii) products of completely monotone functions are completely monotone.
Prove (ii) by Leibniz's rule. Why is this not enough to put Yukawa crystals at Tier~A in this book?
\end{exercise}

\begin{exercise}\label{ex:ch4-image}
Derive the image-charge interaction for two electrons at height $d$ above a grounded plane, and its expansion
$\frac{2d^2}{r^3}$ (in units of $e^2/4\pi\varepsilon_0$) for $r\gg d$. Compute the next term. Which Riesz exponent
governs the large-distance behaviour, and is it in the finite-energy class?
\end{exercise}

\begin{exercise}\label{ex:ch4-density}
For the Riesz energy per particle $e_s(\rho)=c_s\rho^{s/2}$, show that a relative density error $\delta\rho/\rho$
gives a relative energy error $\frac{s}{2}\frac{\delta\rho}{\rho}$ to first order. Evaluate it for $s=3$ and
explain why this matters for the L1 kill criterion.
\end{exercise}

\begin{exercise}\label{ex:ch4-stiff}
For the gradient flow $\dot x=-\mathcal H x$ with $\mathcal H$ symmetric positive semidefinite, show that forward
Euler is stable iff $h<2/\lambda_{\max}$, and that backward Euler is stable for all $h>0$. Estimate the number of
forward-Euler steps needed to reduce the softest non-zero mode by a factor $e$. Why do the zero (translation) modes
not matter?
\end{exercise}

\begin{exercise}\label{ex:ch4-code}
Modify the point-vortex script: (a) replace DOP853 by RK45 and tighten or loosen \texttt{rtol}; plot the relative
drift of $H$ against \texttt{rtol}. (b) Give SciPy's BDF the analytic Jacobian of \eqref{eq:ch4-kirchhoff} and
compare \texttt{nfev}. (c) Change the seed until two opposite vortices come very close, and explain the effect on
the drift of $H$. Report your actual output.
\end{exercise}

\begin{exercise}\label{ex:ch4-riesz3}
Compute, to six digits, the lattice sums $\sum_{a\in\Lat\setminus0}|a|^{-3}$ for the triangular lattice and
$\sum_{a\in\Z^2\setminus0}|a|^{-3}$ for the square lattice, both of covolume one. (Sum over a large disk and correct
the tail by an integral; test your tail correction by doubling the radius.) Which is smaller, and what does the
theorem predict? Use \lean{norm_sq_triangularPoint} for the triangular norms.
\end{exercise}
